\section*{Розділ IV. Особливості проведення різних типів виборів}
\addcontentsline{toc}{section}{Розділ IV. Особливості проведення різних типів виборів}

\subsection*{4.1. Вибори делегатів до Конференції студентів Університету}
\addcontentsline{toc}{subsection}{4.1. Вибори делегатів до Конференції студентів Університету}
    4.1.1. Вибори делегатів до КСУ проводяться за мажоритарною багатомандатною системою з фіксованими квотами від структурних підрозділів Університету.

    4.1.2. Кількість делегатів від кожного факультету/інституту визначається пропорційно до чисельності студентів за формулою: один делегат на кожні розпочаті 500 студентів, але не менше одного делегата від факультету/інституту. Чисельність студентів для розрахунку квоти визначається станом на 1 (перше) число місяця оголошення виборів за офіційними списками адміністрації Університету/ЄДЕБО.

    4.1.3. Додатково до представництва від факультетів/інститутів обираються делегати від студентів, що проживають у гуртожитках, за формулою: один делегат на кожні розпочаті 200 мешканців гуртожитку, але не менше одного делегата від гуртожитку та не більше 5 делегатів від одного гуртожитку. Чисельність мешканців для розрахунку квоти визначається станом на 1 (перше) число місяця оголошення виборів за офіційними списками адміністрації Університету/ЄДЕБО.

    4.1.4. Виборчими округами для виборів делегатів КСУ є:
        \begin{enumerate}[label=\alph*)]
            \item Факультети/інститути (кожен окремо);
            \item Гуртожитки (кожен окремо);
        \end{enumerate}

    4.1.5. Кандидати в делегати висуваються шляхом самовисування.

    4.1.6. Голосування проводиться за окремих кандидатів. У бюлетені виборець ставить позначку "За" напроти тих кандидатів, яких підтримує. Кількість позначок не може перевищувати кількість мандатів делегатів КСУ від відповідного виборчого округу.

    4.1.7. Обраними делегатами вважаються кандидати, які отримали найбільшу кількість голосів у межах квоти відповідного виборчого округу.




\subsection*{4.2. Особливості документообігу під час перших виборів}
\addcontentsline{toc}{subsection}{4.2. Особливості документообігу під час перших виборів}
    4.2.1. Усі документи, пов'язані з першими виборами делегатів КСУ, ведуться ТВК в електронному форматі з обов'язковим створенням резервних копій.

    4.2.2. Протоколи засідань ТВК, рішення про реєстрацію кандидатів, результати голосувань та інші офіційні документи підписуються кваліфікованими електронними підписами уповноважених осіб.

    4.2.3. Вся документація перших виборів зберігається ТВК до моменту передачі архіву першому складу КСУ, після чого передається на зберігання КСУ на строк не менше 3 років.

    4.2.4. Копії всіх документів перших виборів також передаються для архівного зберігання як частина нормативно-правової бази студентського самоврядування.